\chapter{Milestone M1.3: Spectral Projection, Correlator Integration, and DIIS Acceleration}
\label{ch:m1_3}

\section{Zero-Temperature Spectral Projector}
The reciprocal two-point correlation matrix $P(\mathbf{k}) = 2\Pi^-(\mathbf{k})$ projects onto the occupied negative-energy subspace:
\begin{equation}
P(\mathbf{k}) = 2 \sum_{\mathcal{E}_n(\mathbf{k}) < 0} \mathbf{w}_n(\mathbf{k}) \mathbf{w}_n^\dagger(\mathbf{k}) + \sum_{\mathcal{E}_n(\mathbf{k}) = 0} \mathbf{w}_n(\mathbf{k}) \mathbf{w}_n^\dagger(\mathbf{k}).
\end{equation}
The half-weight assignment on zero modes preserves particle-hole symmetry $P(-\mathbf{k}) = 2\mathbb{1} - P^T(\mathbf{k})$.

Integrating $P(\mathbf{k})$ over the Brillouin zone yields real-space covariances:
\begin{align}
\mathrm{U}_b^{\mu\nu} &= \frac{1}{N_{\mathbf{k}}} \sum_{\mathbf{k}} \left( \operatorname{Im}[P_{s_i\mu, s_j\nu}(\mathbf{k})] \cos\phi_b(\mathbf{k}) - \operatorname{Re}[P_{s_i\mu, s_j\nu}(\mathbf{k})] \sin\phi_b(\mathbf{k}) \right),\\
\mathrm{V}_s^{\mu\nu} &= \frac{1}{N_{\mathbf{k}}} \sum_{\mathbf{k}} \operatorname{Im}[P_{s\mu, s\nu}(\mathbf{k})].
\end{align}

\section{The K-006 Invariant and Pearson Random-Walk Average}
In the pure Kitaev ground state ($K = -1.0, u_{\mathrm{gauge}} = +1.0$), the self-consistent matter bond correlator evaluates to:
\begin{equation}
U_b^{00} = -\bar{U} = -\frac{1}{3} \int_0^1 dk_1 \int_0^1 dk_2 |1 + e^{-2\pi i k_1} + e^{-2\pi i k_2}|.
\end{equation}
Remarkably, $3\bar{U} = \mathbb{E}[|\mathbf{R}_3|] = W_3(1)$ is the expected distance from the origin of a 3-step uniform planar Pearson random walk. It possesses the exact closed form:
\begin{equation}
W_3(1) = \frac{3 \cdot 2^{1/3}}{16\pi^4}\Gamma\left(\tfrac{1}{3}\right)^6 + \frac{27 \cdot 2^{2/3}}{4\pi^4}\Gamma\left(\tfrac{2}{3}\right)^6 \approx 1.5745972375506,
\end{equation}
giving $\bar{U} \approx 0.5248657458506$. Midpoint quadrature on meshes up to $12288 \times 12288$ confirms cubic convergence ($p = 3.0000$).

\section{Joint Anderson Acceleration (DIIS)}
Standard Picard updates oscillate due to the hyperbolic saddle-point geometry. Joint DIIS tracks the concatenated vector $X = (\operatorname{vec}(\mathrm{U}), \operatorname{vec}(\mathrm{V}), \boldsymbol{\lambda})^T$ and minimizes historical residuals:
\begin{equation}
\begin{pmatrix} B & -\mathbf{1} \\ -\mathbf{1}^T & 0 \end{pmatrix} \begin{pmatrix} \mathbf{c} \\ \mu \end{pmatrix} = \begin{pmatrix} \mathbf{0} \\ -1 \end{pmatrix}, \qquad B_{pq} = \langle R_{k-p}, R_{k-q} \rangle + \epsilon_{\mathrm{reg}} \delta_{pq}.
\end{equation}
This dampens matter-gauge cross-talk and restores asymptotic linear contraction.
